\chapter{CONCLUSIONS AND FUTURE WORK}

This report presents a measurement-first experiment platform for comparing SLM--LLM collaborative architectures and reports the results of 109 experimental runs executed across five public benchmarks, covering every architecture defined in the Chapter~3 taxonomy. This chapter synthesizes the principal findings, revisits the four research questions from Chapter~1, identifies the limitations of the current evaluation, and sets out the work that remains before the final thesis submission.

\section{Principal Findings}

\paragraph{Platform maturity.}
The experiment platform has moved well beyond a proof-of-concept state. A shared architecture contract---in which the same \texttt{Query} object is processed by any architecture and returns the same \texttt{Response}---makes monolithic, hybrid, ensemble, and swarm-style strategies directly comparable under identical benchmark conditions. The instrumentation layer records latency, token counts, cost, energy, and estimated CO$_2$ for every query. MLflow-backed tracking and a JSON/Markdown reporting pipeline make results reproducible and auditable. This shared measurement policy was the primary platform contribution targeted for this term report, and it is now operational.

\paragraph{Efficiency--quality frontier.}
The most consistent finding across the benchmark suite is that routing and speculative decoding using the \texttt{gemma4-4b}$\rightarrow$\texttt{gemma4-31b} pair occupy the top positions on the efficiency--quality frontier on ARC, HellaSwag, and TruthfulQA. On these benchmarks, routing EATS values of 0.884, 0.869, and 0.844 represent improvements of 15--27 EATS points over the best standalone options, achieved at sub-2-second algorithmic latency and per-query costs below \$0.001. Speculative decoding reaches 1--5 percentage points higher accuracy than routing on these benchmarks at EATS values of 0.817--0.856. This pattern directly addresses the central thesis hypothesis: well-designed collaborative architectures can match or approach standalone LLM accuracy while delivering substantial efficiency gains.

\paragraph{Standalone MoE checkpoints as a competitive baseline.}
Among standalone models, the mixture-of-experts checkpoints \texttt{qwen3.5-35b-a3b} and \texttt{gemma4-26b-a4b} are consistently the most efficient options. They achieve EATS values of 0.492--0.738 across benchmarks---generally above the dense-model standalone range---because their sparse activation profiles keep latency and cost low relative to their accuracy. Any thesis conclusion about collaborative efficiency gains must account for this: the relevant comparison is not against an expensive dense LLM but against the best-efficiency standalone alternative.

\paragraph{Swarm and ensemble accuracy ceilings.}
Pure Swarm achieves the highest raw accuracy on ARC (96.0\%) using no heavy-model calls, and the four-SLM ensemble reaches 96.0\% on GSM8K---within half a point of the best standalone checkpoint---demonstrating that compact peer coordination can reach frontier-level accuracy on structured reasoning tasks. Both findings corroborate multi-SLM collaboration results from the companion systematic literature review. However, both architectures are penalized heavily on EATS by their high wall-clock latency, which reflects the current gap between architectural capability and deployment-level efficiency. Reducing multi-agent coordination latency is therefore as important an engineering challenge as improving accuracy.

\paragraph{Active Oracle and the cost of oracle queries.}
Active Oracle consistently underperforms both routing and speculative decoding in accuracy and EATS despite using the same SLM--LLM pair. The sub-query mechanism incurs additional latency and cost without reliably recovering the correct answer on multiple-choice tasks, indicating that targeted oracle consultation is a poor fit for benchmark formats where the answer is a single letter rather than a free-form explanation. Active Oracle may show different behavior on open-ended generation tasks; the current evidence is specific to the multiple-choice evaluation surface.

\paragraph{Decentralised blackboard coordination.}
The newly benchmarked Blackboard and Entropy Blackboard architectures demonstrate that bid-based decentralised task allocation with an LLM sweeper fallback is empirically viable: on ARC and GSM8K they reach 92.0--97.0\% accuracy while resolving 72--93\% of queries on SLM workers alone, and the 97.0\% on GSM8K---reached by both the highest-bid Blackboard and the first-threshold Entropy Blackboard, the latter with only 7\% sweeper escalation---is the highest raw accuracy in the study. Every reported configuration is a single 100-sample run at the principal operating point, so all design--policy comparisons are setting-matched, and the 2$\times$2 matrix is complete for both architectures. Across that matrix the claim policy proves to be a minor knob---standard Blackboard varies by at most three points between policies and Entropy Blackboard by at most two---and the dramatic policy swings seen in an earlier snapshot (notably 83.0\% vs.\ 69.0\% on TruthfulQA) were artefacts of a legacy short time-to-live rather than the claim rule, with both matched policies landing at 68.0--69.0\% there. What actually moves accuracy is the entropy penalty: it recovers three to seven points on the poorly calibrated benchmarks (HellaSwag, MMLU, TruthfulQA) by escalating the queries overconfident workers would otherwise claim. A preliminary bid-threshold exploration on MMLU shows a monotone escalation--accuracy response (74.0\% to 79.0\% as the threshold rises from 0.65 to 0.85), identifying the bidding rule as the architecture's principal tuning surface. LLM-arbitrated Debate complements this picture, outperforming the task-board designs on factual calibration (87.0\% on TruthfulQA) at a higher fixed orchestration cost.

\paragraph{Measurement lessons.}
One methodological issue identified in the earlier evaluation snapshot was resolved during this reporting period. The GSM8K anomaly for GPT-family models (initially 67.8\% and 59.4\% accuracy) was confirmed to be an output-parsing artifact: after correcting the numeric answer-extraction logic and rescoring the affected runs from stored raw responses, both checkpoints score in line with their performance on the rest of the suite (96.2\% and 94.8\%). The episode demonstrates that free-form benchmarks couple capability measurement to harness-side parsing, and per-model-family validation of the extraction path is now part of the evaluation protocol.

\section{Research Questions Revisited}

\textbf{RQ1 (Performance):} Routing and speculative decoding match or approach the accuracy of the best standalone models on ARC, HellaSwag, and TruthfulQA. Pure Swarm exceeds every standalone model on ARC. The present results provide support for an affirmative answer to RQ1 on these benchmarks: collaborative architectures can match standalone LLM accuracy. On MMLU, Pure Swarm reaches the highest reported accuracy, while routing and speculative decoding remain competitive but do not exceed the top standalone or swarm configurations.

\textbf{RQ2 (Efficiency):} On ARC, HellaSwag, and TruthfulQA, routing reduces per-query cost by factors of 3--12 and latency by factors of 2--7 relative to the same LLM running in standalone mode, while maintaining near-monolithic accuracy. The MoE standalone checkpoints represent the most efficient non-collaborative alternative. Ensemble and Pure Swarm architectures achieve competitive accuracy at the cost of elevated latency. The efficiency picture is therefore architecture- and workload-dependent rather than universal.

\textbf{RQ3 (Unified Metric):} EATS yields a consistent ranking signal across the evaluated benchmarks, capturing the joint accuracy--efficiency trade-off in a single composite score. On ARC, HellaSwag, and TruthfulQA, EATS correctly identifies routing and speculative decoding as the top-ranked configurations and separates the efficient MoE standalone checkpoints from the less efficient dense-model alternatives.

\textbf{RQ4 (Empirical Grounding):} The platform-measured results corroborate several findings from the companion systematic literature review. The ability of SLM-first routing to achieve near-LLM accuracy while limiting heavy-model invocations was a frequently reported trend in the reviewed studies; the ARC, HellaSwag, and TruthfulQA results confirm this empirically under a controlled measurement policy. The ensemble accuracy ceiling on GSM8K and the Pure Swarm result on ARC are consistent with multi-SLM coordination findings from the literature. The Active Oracle underperformance is consistent with reports that oracle-query mechanisms are sensitive to task format. The gain of Entropy Blackboard over standard Blackboard on poorly calibrated tasks empirically supports the emphasis in the uncertainty-aware inference literature on calibrating when a small model should abstain rather than merely measuring its confidence.

\section{Limitations}

The current evaluation has three principal limitations.

First, all five benchmarks use multiple-choice or fixed numeric-answer formats with deterministic gold labels. This makes evaluation straightforward and reproducible but limits coverage of open-ended generation quality, which requires different scoring approaches. The five benchmarks span knowledge breadth, structured reasoning, commonsense inference, arithmetic, and adversarial truthfulness, but they do not yet cover domain-specific or instruction-following tasks.

Second, all hybrid and task-board experiments---routing, speculative decoding, Active Oracle, Debate, Blackboard, and Entropy Blackboard---use a single SLM--LLM pair (\texttt{gemma4-4b}$\rightarrow$\texttt{gemma4-31b}). It is unknown whether the efficiency gains generalize to other pair configurations or whether the specific confidence-threshold setting contributes substantially to the observed results.

Third, the most recently benchmarked architectures---Debate (POA), Blackboard, and Entropy Blackboard---currently rest on a single 100-sample run per configuration, compared with 500 or more samples for the principal standalone baselines, and their orchestration parameters (bid threshold, time-to-live, entropy-penalty weight) have been evaluated at a single principal setting, apart from a preliminary bid-threshold exploration on MMLU. Their results should therefore be read as first empirical evidence rather than final characterisations.

\section{Future Work}

\begin{enumerate}
    \item \textbf{Blackboard-family parameter sensitivity.} With the design--policy matrix now complete and a preliminary bid-threshold exploration on MMLU showing a monotone escalation--accuracy response, systematic sweeps over the bid threshold, the time-to-live before sweeper fallback, and the entropy-penalty weight---across all five benchmarks rather than MMLU alone---will characterise how escalation behaviour, and therefore the accuracy--efficiency balance, responds to these parameters, and will deepen the sample counts for the newly benchmarked architectures toward parity with the standalone baselines.

    \item \textbf{Routing-threshold and pair ablations.} The confidence threshold used for routing (\texttt{gemma4-4b}$\rightarrow$\texttt{gemma4-31b}) has not been systematically varied. Ablations across threshold values and across alternative SLM--LLM pairs will clarify whether the observed EATS gains are robust to configuration choice or specific to the current setting.

    \item \textbf{Human preference evaluation.} For open-ended tasks where exact-match correctness does not fully capture response quality, a UI-backed human preference workflow is planned as a complementary evaluation surface. This component is defined in the repository but not yet operational.

    \item \textbf{Expanded sample counts and experiment diversity.} Reported per-configuration sample counts in the current benchmark suite range from 100 to 520. Increasing sample counts to a consistent minimum of 500 per configuration across all benchmarks would reduce variance in accuracy and EATS estimates and strengthen the statistical reliability of architecture comparisons. Experiment diversity should also be extended to cover a broader set of SLM candidates, additional routing-threshold values, and larger model-tier combinations not yet represented in the registry.

    \item \textbf{Edge device deployment study.} If a collaborative architecture demonstrates a strong EATS profile---specifically, high accuracy at low per-query cost and latency---its feasibility on resource-constrained edge hardware should be investigated. Running a quantized SLM locally on an edge device for the draft stage, with selective escalation to a remote LLM only when the SLM confidence is insufficient, would extend the platform's deployment surface beyond cloud GPU hosts and open a new efficiency dimension not captured by the current cost model.
\end{enumerate}
